\begin{table}[H]
\centering
\caption{Standardized Effect Sizes for Main Outcomes}
\label{tab:sde}
\begin{threeparttable}
\footnotesize
\begin{tabular}{@{}llcccccc@{}}
\toprule
Outcome & Specification & $\hat{\beta}$ & SE & SD($Y$) & SDE & SE(SDE) & Class. \\
\midrule
\multicolumn{8}{@{}l}{\textit{Panel A: Pooled}} \\
Drug OD & Oil $\times$ Boom & 0.142 & 0.226 & 3.77 & 0.038 & 0.060 & Small pos. \\
Drug OD & Oil $\times$ Bust & 0.084 & 0.404 & 3.77 & 0.022 & 0.107 & Small pos. \\
\midrule
\multicolumn{8}{@{}l}{\textit{Panel B: Heterogeneous (High Pre-Boom Drug Rate)}} \\
Drug OD & Oil $\times$ Boom $\times$ High & $-$0.719 & 0.251 & 3.34 & $-$0.215 & 0.075 & Large neg. \\
Drug OD & Oil $\times$ Bust $\times$ High & $-$1.322 & 0.414 & 3.34 & $-$0.396 & 0.124 & Large neg. \\
\bottomrule
\end{tabular}
\begin{tablenotes}[flushleft]
\small
\item \textit{Notes:} \textbf{Country:} United States. \textbf{Research question:} Whether county-level exposure to oil and gas extraction during the shale boom (2005--2014) and subsequent bust (2015) affected drug overdose mortality rates, and whether this effect differed by pre-existing vulnerability to drug deaths. \textbf{Policy mechanism:} The shale revolution created geographically concentrated employment and income shocks in counties overlying shale formations; the mechanism is economic opportunity as a protective factor against substance abuse, operating through labor demand, wages, and community investment in resource-dependent areas. \textbf{Outcome definition:} Model-based age-adjusted drug overdose death rate per 100,000 population from CDC NCHS, covering all drug poisoning deaths (ICD-10 X40--X44, X60--X64, X85, Y10--Y14). \textbf{Treatment:} Binary indicator for county having at least one NAICS 211 (Oil and Gas Extraction) establishment in the pre-boom period (2001--2004), interacted with boom/bust period indicators. \textbf{Data:} CDC NCHS Drug Poisoning Mortality by County, 1999--2015; Census County Business Patterns 2001--2004; county-year panel with 53,387 observations across 3,141 counties. \textbf{Method:} Two-way fixed effects (county + year) difference-in-differences with state-clustered standard errors; triple-difference specification interacting oil exposure with pre-boom drug vulnerability. \textbf{Sample:} All US counties with non-missing drug overdose rates in the CDC NCHS data, 1999--2015. SDE $= \hat{\beta} / \text{SD}(Y)$ where SD($Y$) is the pre-treatment (1999--2004) standard deviation. Classification refers to magnitude, not statistical significance: Large ($|$SDE$| > 0.15$), Moderate ($0.05$--$0.15$), Small ($0.005$--$0.05$), Null ($< 0.005$).
\end{tablenotes}
\end{threeparttable}
\end{table}
